The pilot ran on a preliminary frame built from a leading portion of the bulk
opinion table, before the full pass was committed to. Its purpose was to find
out whether the guideline could be applied at all, whether the controls behaved,
and which strata were worth sampling heavily. Sixty items were annotated. Pilot
items are not part of the benchmark and no pilot label appears in any result.

\subsection{What the pilot showed} Labels fell out by stratum as follows:
\textsc{ctrl\_hold} 12 \textsc{decided} and 0 \textsc{unresolved};
\textsc{plain} 19 \textsc{unresolved} and 1 \textsc{decided};
H2 8 \textsc{unresolved} and 2 \textsc{decided};
H1 5 \textsc{unresolved}, 5 \textsc{decided}, and 5 \textsc{unclear}.

Three things follow. The holding-cue control behaves as intended and is not
contaminated. The \textsc{plain} stratum is nearly degenerate: a trigger with no
structural complication almost always does mean the court left the issue open,
which is why a keyword rule scores as well as it does and why that score means
little. The H1 stratum is close to an even split, so it is where the task
actually lives, and the sample was reallocated toward it. This reallocation
changes only how many items are drawn per stratum. Label definitions and stratum
definitions were fixed before the pilot and held throughout.

\subsection{The one protocol revision} Version 1.0 of the guideline had no rule
for a passage whose only trigger sits inside a parenthetical or a description of
another court's work, on a question the present opinion never reaches in its own
right. A citation noting that the Supreme Court declined to decide whether an
airport user fee was excessive tells you nothing about the disposition of the
Federal Circuit opinion that cites it, and forcing a choice between
\textsc{decided} and \textsc{unresolved} would have manufactured a label. Rule
L1-9 routes these to \textsc{unclear}. It uses an existing label; it adds none.
This is the single revision the design permits, it was made once, before any
main-sample annotation, and the changelog is reproduced in
Appendix~\ref{app:guide}.

\subsection{Three pipeline defects the pilot exposed} These concern candidate
construction, not the guideline. The pattern flagging a trigger as plausibly
describing another institution was case-sensitive, matching ``the district
court'' but missing ``the District Court,'' which appellate opinions write about
as often; it now also recognizes ``the court held'' and its near neighbours. The
annotator was shown a later holding sentence only for items already flagged H2,
though Rule L1-3 turns on precisely that passage, so an H1 item whose opinion
later resolved the issue was unanswerable rather than hard; the sentence is now
shown whenever one exists outside the displayed window. And ``our holding'' was
a control locator, but in the pilot it pointed at a named prior case rather than
the present disposition every time it fired, so it was dropped.

\subsection{What the pilot did not do} It was not used to adjust any label
boundary. The trigger dictionary was frozen before the pilot and is unchanged;
in particular, the pilot surfaced two recurring dictionary false positives, and
neither was patched. One is the statutory-scope reading of \emph{reach}, as in
``the shipment clause does not reach shipments between the continental States,''
which is a holding about a statute rather than a non-resolution. The other is a
line-break artefact in the source text that splits \emph{preserving} into
\texttt{P} and \texttt{reserving}, which then matches T11. Both stay in, are
labelled \textsc{decided} by the guideline, and are counted against the
dictionary in \S\ref{sec:models}.
